% BÀI TOÁN HÌNH HỌC - BAI_HINH_27 (TRÍCH ĐỀ 27)
\begin{ex}
Cho đường tròn (O), bán kính R ($R>0$) và dây cung BC cố định. Một điểm A chuyển động trên cung lớn BC sao cho tam giác ABC có ba góc nhọn. Kẻ các đường cao AD, BE của tam giác ABC cắt nhau tại H và BE cắt đường tròn (O) tại F (F khác B).
	\begin{enumerate}
		\item[1)] Chứng minh rằng tứ giác DHEC nội tiếp.
		\item[2)] 
		\begin{enumerate}
			\item[a)] Kẻ đường kính AM của đường tròn (O) và OI vuông góc với BC tại I. Chứng minh rằng I là trung điểm của HM.
			\item[b)] Khi BC cố định, xác định vị trí của A trên đường tròn (O) để $DH \cdot DA$ lớn nhất.
		\end{enumerate}
	\end{enumerate}
\end{ex}
